\section{Conclusion}
\label{sec:conclusion}

This review connects scientific-object classification to agent decisions, the evidence those decisions generate, and the results available for subsequent research. Its six questions lead to the following conclusions.

\textbf{RQ1: Domains and objects.} The 660 domain-assigned works span the eleven modules, alongside two works retained in the separate general-method bucket. Non-exclusive assignments map evaluated objects and tasks. The distribution describes the assembled corpus, with substantial variation within each broad domain.

\textbf{RQ2: Research roles.} Agents connect information acquisition, hypothesis development, design, execution, analysis, validation, and communication in different combinations. Their shared contribution is organizing connected actions. Scientific interpretation requires identifying the consequential choices within those activities, rather than counting the stages covered.

\textbf{RQ3: Decision paradigms.} Symbolic reasoning, model-based planning, probabilistic optimization, learned policies, LLM reasoning, and hybrid arrangements implement complementary decisions. Matched comparisons locate gains in information use, candidate selection, generated procedures, or tool access. Architectural complexity alone does not identify their source or establish a universal ordering of paradigms.

\textbf{RQ4: Autonomy and collaboration.} The cases demonstrate distinct allocations of candidate selection, hypothesis discrimination, calibration, observation scheduling, and explanatory-model revision. Researcher-supplied representations and human experimental work remain consequential. The scientific effect of delegation depends on which assumptions, measurements, or subsequent actions the system can change in response to evidence.

\textbf{RQ5: Scientific evidence.} Feedback advances investigation by informing decisions relevant to the claim or revealing that its current formulation is inadequate. Improving candidate performance, predicting a response throughout a domain, and discriminating explanations are different objectives. Adaptive acquisition can help the first while weakening the second or leaving the third unresolved. The cases establish distinct scientific products: a checked witness supports a mathematical bound, an executable heuristic can generate further solutions, controlled follow-up explains newly exposed behavior, and an experimentally tested rule guides later investigation. These products require different justification. A proposed discriminating prediction remains a useful but uncompleted empirical test. Performance attribution and the scientific content of the output must therefore be assessed separately.

\textbf{RQ6: Cross-domain comparison.} Disciplinary conditions shape the appropriate research strategy. Explicit verification can separate a mathematical search from checking its result; variable physical responses call for measurement and reproducibility assessment. Transient phenomena place value on timely acquisition, while proxy endpoints require additional evidence when the intended claim concerns a more distant outcome. Transfer can preserve a procedure without preserving its comparative advantage or the scientific relation inferred from its output. These conditional findings explain participation more precisely than disciplinary maturity or data volume alone, while leaving causal explanations of publication counts unresolved.

The focused model--observation comparison gives these conclusions physical content: learning a predictive discrepancy, reconstructing a measurement-to-parameter relation, and changing an explanatory model address different uncertainties. The relevant achievement is a justified choice of what to investigate next and a result whose intended use is supported by the evidence. This provides a common basis for evaluating scientific agents across domains while retaining the differences that make their scientific claims meaningful.
